\subsection{幂零李代数的定义}

定义 1. 若李代数 \(\mathfrak{g}\) 有一个理想的有限递降序列 \((\mathfrak{g}_i)_{1\leqslant i\leqslant p}\)，其中 \(\mathfrak{g}\)，且对 \(\mathfrak{g}_0 = \mathfrak{g},\mathfrak{g}_p = \{0\}\) 有 \([\mathfrak{g},\mathfrak{g}_i]\subset \mathfrak{g}_{i + 1}\)，则称 \(0\leqslant i < p\) 是幂零的。

交换李代数是幂零的。

命题 1. 设 \(\mathfrak{g}\) 是李代数。以下条件等价：

\begin{enumerate}
\def\labelenumi{(\alph{enumi})}
\item
  g 是幂零的；
\item
  对充分大的 \(\mathcal{C}^k\mathfrak{g} = \{0\}\)，\(k\)；
\item
  对充分大的 \(\mathcal{C}_k\mathfrak{g} = \mathfrak{g}\)，\(k\)；
\item
  存在整数 \(k\)，使得对 \(\operatorname{ad} x_1 \circ \operatorname{ad} x_2 \circ \dots \circ \operatorname{ad} x_k = 0\) 中的所有元素 \(x_1, x_2, \ldots, x_k\)，都有 \(\mathfrak{g}\)；
\item
  存在 \((\mathfrak{g}_i)_{0\leqslant i\leqslant n}\) 的理想递降序列 \(\mathfrak{g}\)，其中 \(\mathfrak{g}_0 = \mathfrak{g},\mathfrak{g}_n = \{0\}\)，且对 \([\mathfrak{g},\mathfrak{g}_i]\subset \mathfrak{g}_{i + 1}\) 有 \(\dim \mathfrak{g}_i / \mathfrak{g}_{i + 1} = 1\) 以及 \(0\leqslant i <   n\)
\end{enumerate}

若 \(\mathcal{C}^k\mathfrak{g} = \{0\}\)（分别为 \(\mathcal{C}_k\mathfrak{g} = \mathfrak{g}\)），显然序列 \(\mathcal{C}^1\mathfrak{g},\ldots ,\mathcal{C}^k\mathfrak{g}\)（分别为 \(\mathcal{C}_k\mathfrak{g},\mathcal{C}_{k - 1}\mathfrak{g},\ldots ,\mathcal{C}_0\mathfrak{g}\)）具有定义 1 的性质，因而 \(\mathfrak{g}\) 是幂零的。反之，假设存在具有定义 1 所述性质的序列 \((\mathfrak{g}_i)_{0\leqslant i\leqslant p}\)。对 \(i\) 作归纳可见，\(\mathfrak{g}_i\supset \mathcal{C}^{i + 1}\mathfrak{g}\) 且 \(\mathfrak{g}_{p - i}\subset \mathcal{C}_i\mathfrak{g}\)。因此 \(\mathcal{C}^{p + 1}\mathfrak{g} = \{0\}\) 且 \(\mathcal{C}_p\mathfrak{g} = \mathfrak{g}\)。至此证明了条件 (a)、(b)、(c) 等价。另一方面，\(\mathcal{C}^i\mathfrak{g}\) 是形如

\[
\left[ x _ {1}, \left[ x _ {2}, \dots , \left[ x _ {i - 2}, \left[ x _ {i - 1}, x _ {i} \right] \right] \dots \right] \right]
\]

其中 \(x_{1}, x_{2}, \ldots, x_{i}\) 遍历 \(\mathfrak{g}\)。因此条件 (b) 和 (d) 等价。最后，若存在具有定义 1 所述性质的理想序列 \((\mathfrak{g}_i)_{0 \leqslant i \leqslant p}\)，

定义 1，则存在一个维数为 \((\mathfrak{h}_i)_{0\leqslant i\leqslant n}\) 的 \(\mathfrak{g}\) 的向量子空间递降序列 \(n,n - 1,n - 2,\ldots ,0\)，以及指标序列

\[
i _ {0} <   i _ {1} <   \dots <   i _ {p}
\]

其中 \(\mathfrak{g}_0 = \mathfrak{h}_{i_0},\mathfrak{g}_1 = \mathfrak{h}_{i_1}\dots ,\mathfrak{g}_p = \mathfrak{h}_{i_p}\)；由于 \([\mathfrak{g},\mathfrak{h}_{i_k}]\subset \mathfrak{h}_{i_{k + 1}}\)，\(\mathfrak{h}_i\) 是理想，并且对所有 \([\mathfrak{g},\mathfrak{h}_i]\subset \mathfrak{h}_{i + 1}\) 都有 \(i\)。因此条件 (a) 和 (e) 等价。

推论 1. 非零幂零李代数的中心非零。

推论 2. 幂零李代数的 Killing 形式为零。

在幂零李代数中，对所有 \(x\) 和 \(y\)，ad \(x \circ\) ad \(y\) 都是幂零的，因而迹为零。

命题 2. 幂零李代数的子代数、商代数和中心扩张都是幂零的。有限个幂零李代数的直积也是幂零李代数。

设 g 是李代数，g' 是 g 的子代数，b 是 g 的理想，t = g/b，φ 是从 g 到 t 的典范映射。若 g 幂零，则对某个整数 k 有 \(C^{k}g = \{0\}\)，从而 \(C^{k}g' \subset C^{k}g = \{0\}\)，且 \(C^{k}t = \phi(C^{k}g) = \{0\}\)，所以 g' 和 t 都幂零。若 t 幂零且 b 包含于 g 的中心，则对某个整数 k 有 \(C^{k}t = \{0\}\)，从而 \(C^{k}g \subset b\)，于是 \(C^{k+1}g \subset [b, g] = \{0\}\)，故 g 幂零。最后，关于直积的断言例如由命题 1 的断言 (a) \(\Leftrightarrow\) (d) 得到。

定义 1 和命题 2 表明，幂零李代数恰好是通过一系列中心扩张从交换李代数得到的代数。

命题 3. 设 g 是幂零李代数，h 是 g 的一个不等于 g 的子代数。h 在 g 中的正规化子不等于 h。

设 \(k\) 是满足 \(\mathcal{C}^k\mathfrak{g} + \mathfrak{h}\neq \mathfrak{h}\) 的最大整数。则

\[
[ \mathcal {C} ^ {k} \mathfrak {g} + \mathfrak {h}, \mathfrak {h} ] \subset \mathcal {C} ^ {k + 1} \mathfrak {g} + \mathfrak {h} \subset \mathfrak {h}
\]

因此 \(\mathfrak{h}\) 在 \(\mathfrak{g}\) 中的正规化子包含 \(\mathcal{C}^k\mathfrak{g} + \mathfrak{h}\)。

